\documentclass[12pt, a4paper]{article}
\usepackage[T2A]{fontenc}
\usepackage[russian]{babel}
\usepackage[margin=0.6in]{geometry}
\usepackage{amssymb}
\usepackage{amsmath}
\usepackage{graphicx}
\usepackage{blindtext}
\usepackage[none]{hyphenat}
\usepackage[expansion=false]{microtype}
\emergencystretch=3em
\def\D{\mathbb}
\graphicspath{{images/}}

\title{Математическа екология}
\author{Матей Захариев}
\date{}
\begin{document}
\maketitle

{\centering\section*{14.10.26}}
\subsection*{Екология, Биосфера, Основни понятия}
Екология - Ecology. Произлиза от oikos (дом) и logia(наука). Основана края на 19-ти век.\\

\textbf{Биосфера}: Въведено 1926, основополагащо на ак. Вернадски. Животът е геологическата сила, формираща Земята.\\
\subsubsection{Различни деления на елементите на биосферата.}
Деление на Вернадски:
\begin{enumerate}
	\item Литосфера: Най-външната твърда обвивка на земята с дебелина между 10 и 50 километра. Основно място за развитие на живот. Най-важен е почвения слой. Геохимичната лаборатория на Биосферата. При неговото присъствие почвения слой е с дебелина между 1 и 2 метра. Присъствие на органични полезни изкопаеми: Нефт, природен газ, въглища.
	\item Хидросфера: Водната обвивка на земята. Друга основна среда за развитието на живота. Водата е най-широко разпространената субстанция в биосгерата. Ключова за много химични реакции, обър разтворител, различни агрегатни състояния, електропроводима в качеството си на разтворител, топлоемка, мобилна в планетарни условия.
	\item Атмосфера: Газовата обвивка на земята. Активни биохимически процеси. Фотосинтеза, дишане, окисляване.
	\item Космическа енергия: Слънчева енергия (UV, IR, светлина, високо-енергийни частици). В биосферата слънчевата енергия се превръща в ползотворна.
	\item Биовещество: Главна структурна единица на Биосферата. Влиае върху всички природни процеси и описаните кръгове. Деление на биовеществото:
	\begin{itemize}
		\item Според храненене и дишане:
		\begin{itemize}
			\item Автотрофни: Самохранеши се. Изхранват се чрез неорганични вещества.
			\item Хетеротрофни: Хранят се чрез органични вещества
		\end{itemize}
		\item Според типа на преобразуване на енергия:
		\begin{itemize}
			\item Продуценти: Хомосинтеза и хемосинтеза;
			\item Консументи 1-ви ред: Енергия от продуценти;
			\item Консументи 2-ри ред: Енергия от консументи 1-ви ред;
			\item ...
		\end{itemize}
	\end{itemize}
\end{enumerate}

Деление по териториален принцип:
\begin{itemize}
	\item Континентална Биосфера. $149*10^6 км^2$, 29\%. Различни фактори, определящи условията: Географска ширина, надморска височина, количество валежи. Влажни тропически, Влажни умерени, Полу пустинни и пустинни.
	\item Океанска биосфера. 71\%, $361*10^6 км^2$. Океанската вода: $370*10^9 км^3$ = 98\% от водата на планетата. Останалите два процента са прясна. За сравнение язовир Искър е $0.6 км^3$. Деление на океанската биосфера: Арктични и Екваториално-тропически по ширина и евфотични и дисфотични по дълбочина(наличие на светлина).
\end{itemize}

\subsubsection{Кръговрати на веществата}
Имаме 5 основни кръговрата:
\begin{enumerate}
	\item Енергичен кръговрат: Постъпване и преобразяване на слънчева енергия в биосферата.
	\item Минерален кръговрат: Преместване на продукти чрез вулканични, магматични процеси и изветряване.
	\item Воден кръговрат: Хидрологична и атмосферна циркулация.
	\item Газов кръговрат: Миграция на различни по произход газове - вулканични, атмосферни, биохимични процеси.
	\item Биологичен кръговрат: Органични вещества във всякакви процеси.
\end{enumerate}

В текущите знания има твърде много пропуски относно това как живата природа влиае на неживата. ДОбре изучени са съдържанията на различните химически велементи в състава на биосферата:
\begin{itemize}
	\item Макроелементи (над 0.1\%)
	\begin{itemize}
		\item Водород, Кислород (над 10\%)
		\item Въглерод, Азот, Калции (1-10\%)
		\item Сяра, Фосфор, Силиций, Калий (0.1-1\%)
	\end{itemize}
	\item Микроелементи (под 0.1\%)
	Желязо, Натрий, Хлор, Алуминий, цинк, бром, мед, ванадий, хром, йод.
	\item Ултраелементи
	Кобалт, Бисмуд, сребро, злато, уран, радий, торий.
\end{itemize}

{\centering\section*{21.10.26}}
\subsection*{Вода, Въздух, Атмосфера.}
\subsubsection*{Вода}
Както споменато в предната лекция, само 2\% от водата на планетата е прясна и само 0.1\% от нея е използваема, освен което тя е и неравномерно разпределена. Човекът консумира 4000 км$^3$ вода годишно:
\begin{itemize}
	\item $\frac{3}{4}$ за земеделие.
	\item $\frac{1}{4}$ за промишленост и бит.
\end{itemize}
Около 500 км$^3$ се превръщат в отпадъчни води.\\

Засушаване причинено от човека. Има множество причини: строене на извори, резервоари, изсичане на гори, т.н. Океанът е голям източник на кислород благодарение на планктон и други автотрофи. Той се замърсява от човека с нефтопродукти/пластмаса. Генерирането на прясна вода: обезсоляване-енергийно скъп процес, добив на вода от полярни шапки. 

\subsubsection*{Въздук и атмосфера}

Състои се от 78\% азот ($N_2$), 21\% кислород($O_2$), 0.2\% въглероден диоксид ($CO_2$) и други. Въглеродният диоксид пропуска светлина, но задържа и отблъсква инфрачервена светлина, създавайки парников ефект. Основен генератор на кислород е фотосинтезата и автотрофите. Озоновият ($O_3$) слой защитава от вредни UV лъчи. Подлежи на химично (вредни елементи, киселинни замърсявания) и механично (прахови частици и други микрочастици) замърсяване.

\subsubsection*{Почва}

Органични вещества, минерали, организми. 70\% от тях са неплодородни. Човешко замърсяване благодарение на агресивни процедури: тежко наторяване, ерозия, химично замърсяване с пестициди и промишлен отпадък, асфалтиране, т.н.

\subsubsection*{Боисфера}

Въздействие на човека върху биосферата. Техническата дейност след ВСВ влияе сериозно върху природата. Промишлената дейност произвежда следните отпадъци - Твърди отпадъци - $40*10^9$ тона на година:
\begin{itemize}
	\item{50\%} органични отпадъци.
	\item{5\%} киселинни агенти.
	Пестициди, перилни препарати, тежки метали.
\end{itemize}
Радиоактивно замърсяване: 3 вида:
\begin{itemize}
	\item $\alpha$ радияция. $He^+$ Ниска пробиваемост - спиран от лист хартия, опасна при консумация или приемане в тялото. Висока йонизираща сила.
	\item $\beta$ радияция. $e^-$ средна пробиваемост - спиран от кожа, тънък лист алуминий или пластмаса, опасна при консумация или приемане в тялото, може да изгаря кожа. Средна йонизираща сила.
	\item $\gamma$ радияция. Високочестотни вълни/фотони, много висока пробиваемост - лесно преминава през човешкото тяло, високо опасна поради което. Слаба йонизираща сила.
\end{itemize}

Всеки радиоактивен исотоп има т.нар. период на полуразпад - времето отнемащо половината от радиацията да се разпръсне. Скоростта на разпръсване е пропорционална на масата, съответно радиацията никога не стига нула. Радиоактивни елементи се натрупват по трофическите вериги. 

\subsubsection*{Енергетика и замърсяване}
\begin{itemize}
	\item Атомна. Замърсяване при аварии. Рядко се случват, но при неправилна реакция, могат да водят до екологическа катастрофа.
	\item Нефт и въглища - емисии и изчерпване на кислород.
	\item Алтернативни - Вец, Фец, Приливни и отливни, Вятър. Важно е как се съхранява енергията - батерии, механични и електрически.
	\item Ядрен синтез - някой ден.
\end{itemize}

Шумово замърсяване, добив на изкопаеми, изменение на релефа. След ВСВ добиваме над 80 елемента.

\subsubsection*{Естествени равновесия}
Кръговрат на веществата.
Неорганични в-ва, автотрофи ги превръщат в органични, сапротрофни превръщат обратно в неорганични. Геохимично равновесие.

Енергично равновесие - Постоянно усвояване на космическата енергия в Биосферата.

\subsection*{Емпирични обобщения за биосферата. (Аксиоми на Вернодски)}

\textbf{Дефиниция}: Емпиричните обобщения са система от аксиоми, които описват текущото ниво на познания за биосферата и се явяват основа за последващо развитие на науката.
\begin{enumerate}
	\item Аксиома за хомогенност.
	В продължение на всички геоложки периоди не се е случвало жива материя да се появи от неживата. ``Живо само от живото''
	N(t) - количество биовещество в момент t. Динамиката му се описва $N'(t) = \rho(N)N(t)$. Съответно, ако в даден момент $N(t_0)=0$, то $N'(t_0) = 0$ и съответно нямаме промяна.
	\item Аксиома за ергодичност.
	Не знаем как възниква животът.
	\item Непрекъснатост. 
	Съвременното биов-во е свързано генетически със биов-вото от предните епохи.
	\item Постоянен химичен състав
	\item Постоянна биомаса
	\item Космическа енергия
	\item Автотрофи
\end{enumerate}
Последните 4 са аксиоми за стационарност и 4. и 5. са най-оспорвани.
\end{document}